\par\Needspace{6\baselineskip}\medskip
\hypertarget{native-120-article089}{}\pdfbookmark[2]{89.}{article-089}
\begin{center}89.\end{center}
We now have two tasks. The first is to find the unknown $g$ as conveniently as possible from the transcendental equation 12, which admits no direct solution. The second is to deduce the elements themselves from the angle $g$ once it has been found. Before taking up these tasks, we shall describe a transformation \SourceLeaf{121}that simplifies the calculation of the auxiliary quantity $l$ or $L$ and gives a more elegant form to several formulas that will be developed later.

Specifically, we introduce the auxiliary angle $\omega$ determined by the formula
\[
\BareRoot[4]{\dfrac{r'}{r}}=\operatorname{tang}(45^0+\omega)
\]
and obtain
\[
\begin{aligned}
\BareRoot{\dfrac{r'}{r}}+\BareRoot{\dfrac{r}{r'}}
&=2+\{\operatorname{tang}(45^0+\omega)-\operatorname{cotang}(45^0+\omega)\}^2\\[4pt]
&=2+4\operatorname{tang}2\omega^2;
\end{aligned}
\]
from which we have
\[
\begin{aligned}
l&=\frac{\sin\tfrac12f^2}{\cos f}+\frac{\operatorname{tang}2\omega^2}{\cos f}\mathpunct{\raisebox{.45ex}{,}}\\[6pt]
L&=-\frac{\sin\tfrac12f^2}{\cos f}-\frac{\operatorname{tang}2\omega^2}{\cos f}\mathpunct{\cdot}
\end{aligned}
\]
